% ID: SEL-D0061-OB-E02
% ESTADO: validada. S061.
\textbf{a)} Por la regla del cociente y la cadena,
\[
f'(x)=\frac{e^{3x}\bigl(3(1+x^2)-2x\bigr)}{(1+x^2)^2},
\quad x\in\mathbb R.
\]
\textbf{b)} El argumento del logaritmo es $x+3x^3$ y es positivo sólo para $x>0$. Aplicamos $(\ln u)'=u'/u$:
\[
g'(x)=\frac{1+9x^2}{x(1+3x^2)},\quad x>0.
\]
\textbf{c)} Derivamos $2^{5x}$ mediante la cadena y escribimos $1/x^2=x^{-2}$:
\[
h'(x)=5(\ln2)\,2^{5x}-\frac{2}{x^3},\quad x\ne0.
\]
